\section{Pruebas Realizadas y Resultados}
\label{sec:pruebas_resultados}

La estrategia de Aseguramiento de Calidad (\textit{Quality Assurance} - QA) de \textbf{artesa} se diseñó bajo una política de pruebas automatizadas continuas, ejecutadas directamente sobre el ejecutor nativo \texttt{node:test} de Node.js 22. Este enfoque garantiza una velocidad de ejecución instantánea y elimina la sobrecarga de dependencias pesadas de prueba.

\subsection{Tipología y Cobertura de las Pruebas}

La batería de pruebas se divide en cinco categorías estratégicas:

\subsubsection{1. Pruebas de Contratos de API REST}
Verifican el estricto apego de los controladores al estándar JSend y la correcta validación de esquemas con \textbf{Zod}:
\begin{itemize}
    \item Validación de cargas útiles (\textit{payloads}) en solicitudes \texttt{POST} y \texttt{PUT}.
    \item Retorno de códigos de estado HTTP semánticos (200 OK, 201 Created, 400 Bad Request, 401 Unauthorized, 403 Forbidden, 404 Not Found y 409 Conflict).
    \item Estructura unificada en respuestas de éxito (\texttt{\{status: "success", data: ...\}}) y error (\texttt{\{status: "fail", data: ...\}}).
\end{itemize}

\subsubsection{2. Pruebas de Máquinas de Estados Finitos (FSM)}
Auditan las transiciones de estado de los flujos de trabajo críticos para prevenir operaciones ilegales:
\begin{itemize}
    \item Ciclo de vida de órdenes del POS: comprobación de transiciones válidas (\texttt{BORRADOR} $\rightarrow$ \texttt{CONFIRMADA} $\rightarrow$ \texttt{EN\_PREPARACION} $\rightarrow$ \texttt{LISTA} $\rightarrow$ \texttt{ENTREGADA}) y rechazo inmediato de transiciones retrógradas o inválidas.
    \item Ciclo de caja: imposibilidad de registrar cobros o movimientos sobre turnos en estado \texttt{CERRADO}.
\end{itemize}

\subsubsection{3. Pruebas Aritméticas Anti-Float}
Comprueban que los cálculos financieros no sufran desviaciones por representación binaria de coma flotante (IEEE 754):
\begin{itemize}
    \item Cálculo de precios con IVA incluido (tasa del 13\%) y desglose matemático inverso para facturación de consumidor final y crédito fiscal.
    \item Retenciones de seguridad social (ISSS 3\%, AFP 7.25\%) con aplicación estricta de bases máximas de cotización.
    \item Conciliación de montos en cobros mixtos (la suma de los pagos debe igualar exactamente al total facturado).
\end{itemize}

\subsubsection{4. Pruebas de Autorización y Matriz RBAC}
Validan la seguridad perimetral de cada ruta y componente:
\begin{itemize}
    \item Rechazo de acceso a usuarios no autenticados (HTTP 401).
    \item Bloqueo de operadores con rol activo pero sin el permiso explícito en base de datos (HTTP 403).
    \item Protección del Kárdex: ocultamiento automático de la columna \texttt{costo\_unitario} cuando el usuario posee \texttt{INV\_VER\_KARDEX} pero carece de \texttt{INV\_VER\_COSTOS}.
    \item Protección de seguridad operativa: verificación de que un usuario no pueda auto-desactivarse ni removerse sus propios privilegios administrativos.
\end{itemize}

\subsubsection{5. Pruebas de Contratos de Interfaz (Frontend)}
Verifican la coherencia visual y de datos en la SPA de React:
\begin{itemize}
    \item Renderizado condicional del menú según permisos asignados en el token de sesión.
    \item Sincronización del estado de acordeón en la barra lateral con la ruta activa del navegador.
    \item Verificación de componentes compartidos (tablas con cabecera fija \textit{sticky}, scroll interior y paginación).
\end{itemize}

\subsection{Resultados Obtenidos}

La Tabla~\ref{tab:resultados_pruebas} resume la distribución de las pruebas automatizadas implementadas en el sistema.

\begin{table}[htbp]
\caption{Resumen de la Suite de Pruebas Automatizadas}
\label{tab:resultados_pruebas}
\begin{tabularx}{\columnwidth}{lXXr}
\toprule
\textbf{Suite / Archivo} & \textbf{Dominio Evaluado} & \textbf{Tipo de Prueba} & \textbf{Tests} \\
\midrule
\texttt{pos-orders-contracts} & Órdenes y FSM & Contrato y Lógica & 25 \\
\texttt{inventario-contracts} & Stock, Recetas, Kárdex & Contratos y Permisos & 19 \\
\texttt{roles-permisos-contracts} & RBAC y Usuarios & Seguridad y Regresión & 12 \\
\texttt{turnos-caja-contracts} & Arqueo y Cuadre & Lógica Financiera & 8 \\
\texttt{empresa-contracts} & Sucursales y Parámetros & Contrato de API & 8 \\
\texttt{frontend-pos-menu} & Menú y Categorías POS & Interfaz y Permisos & 12 \\
\texttt{frontend-rbac-menu} & Navegación y Permisos & RBAC en UI & 14 \\
\texttt{frontend-pages-contracts} & Vistas y Rutas Protegidas & Integración SPA & 16 \\
\texttt{reportes-ux-contracts} & Reportes y Métricas & Presentación de Datos & 14 \\
\midrule
\textbf{Total de Pruebas} & \multicolumn{2}{l}{\textbf{Ejecutadas con éxito (100\% pass)}} & \textbf{128} \\
\bottomrule
\end{tabularx}
\end{table}

\textbf{Métricas de Desempeño:} La ejecución completa de la suite de 128 pruebas automatizadas toma un promedio de tan solo \textbf{0.92 segundos} en un entorno de desarrollo estándar, garantizando retroalimentación inmediata durante el proceso de integración continua sin degradar los tiempos de compilación.
